% ==========================================================================
% Chapter 6: Conclusion (target ~350 words, max 1000)
% ==========================================================================

\chapter{Conclusion (565/1000 words)}
\label{ch:conclusion}

This project set out to determine whether a multi-agent orchestration of distinct pre-trained Claude models, grounded in structured biometric features and constrained by schema validation, can produce Hyrox coaching guidance that is consistent, plausible, and useful. At this stage of the project, the honest answer is: plausible, demonstrably, and not yet consistent.

The strongest evidence gathered so far is the Critic ablation in Chapter~\ref{ch:evaluation}: at the true design tiers, the Critic intervened on every single evaluated run, rejecting outright over a third of the time and forcing a substantive revision the rest. Chapter~\ref{ch:implementation}'s live-run figure gives that statistic a face --- a specific, cited set of safety reasons turning an unsafe seven-day plan into a four-session plan with an explicit rest day. Combined with a fully test-first, hexagonally-structured codebase and a documented recovery from a real model-specific reliability failure, this is credible evidence that the orchestration discipline this project set out to demonstrate is not merely feasible but doing measurable work.

Against that, the same evaluation is equally direct about what is not yet achieved. Output consistency sits at 25\%, meaning the same input rarely produces the same plan twice --- a genuine limitation of the current design, not an artefact of a weak substitute model, and one this report treats as a concrete engineering problem rather than a caveat to be minimised. The user study that would validate whether athletes actually trust and act on this guidance has not started. The feature-engineering layer now implements two of the three planned training-load signals --- the Banister TRIMP score and the acute-to-chronic workload ratio --- leaving only pace/heart-rate zones outstanding. And the implemented schema is simpler than the one originally designed, a trade-off made deliberately but one that narrows how much structure downstream components can currently rely on.

The broader theme this project is accumulating evidence for is that model-role specialisation is not a cosmetic cost-saving choice: swapping the Coach and Critic from a Haiku substitution to their true tiers changed the system's \emph{failure mode}, not just its output quality, and the Critic's intervention rate is itself a measurable function of which model plays that role. Whether that generalises beyond the Claude model family, and whether it holds up once real athletes' data replaces a synthetic fixture, are the two questions the remaining project weeks are built to answer.

The originality of the project lies not in any single model or feature but in their combination: an open, auditable, role-specialised multi-model pipeline --- distinct Claude tiers for classification, generation, and independent critique, with schema-validated hand-offs --- applied to hybrid-sport (Hyrox) coaching, a domain no current commercial system serves with a transparent, separately-reasoning safety stage. The concrete further work that follows from this evaluation is fourfold. First, replace the single 16-session synthetic fixture with the planned 50-plus-session multi-athlete labelled set and run the four-week user study, so consistency and usefulness are measured on real data. Second, constrain the Coach's sampling behaviour and add a consistency rule to the Critic, then re-run the stability metric to attack the 25\% figure directly. Third, extend the Critic ablation across a second model family to test whether role-specialisation holds beyond Claude. Fourth, complete the feature-engineering layer with pace/heart-rate zones so the Coach reasons over the full load model the design intends. Each is bounded and specified, and together they define what would turn a demonstrably plausible system into a demonstrably consistent and trustworthy one.
